% SEGMENT OLP-0451-S01
% Part: normal-modal-logic
% Chapter: filtrations
% Section: introduction

\documentclass[../../../include/open-logic-section]{subfiles}

\begin{document}

\olfileid{nml}{fil}{int}

\olsection{அறிமுகம்}

ஒரு தருக்கத்தைப் பற்றிய முக்கியமான கேள்விகளில் ஒன்று, அது
தீர்மானிக்கத்தக்கதா, அதாவது, ``இந்த !!{formula} ஏற்புடையதா?'' என்ற
கேள்விக்கு விடையளிக்கும் பயனுறு செயல்முறை உள்ளதா என்பதாகும். கூற்றுத்
தருக்கம் தீர்மானிக்கத்தக்கது: ஒரு மெய்மை அட்டவணையை அமைப்பதன் மூலம்
!!a{formula} ஒரு மெய்மமா என்பதை நாம் பயனுறச் சோதிக்கலாம்; கொடுக்கப்பட்ட
!!{formula} ஒன்றுக்கான மெய்மை அட்டவணை முடிவுறுவதாகும். ஆனால் ஒரு
வாய்ப்புநிலை வாய்பாடு எல்லா மாதிரிகளிலும் மெய்யா என்பதை வெளிப்படையான
முறையில் சோதிக்க முடியாது; ஏனெனில் அவை முடிவிலா எண்ணிக்கையில் உள்ளன.
கொடுக்கப்பட்ட வாய்பாட்டுக்குப் பொருத்தமான எல்லா முடிவுறு மாதிரிகளையும்
நாம் பட்டியலிடலாம்; ஏனெனில் அந்த !!{formula} இல் உண்மையில் இடம்பெறும்
!!{propositional variable}s க்கு உலகங்களின் உட்கணங்களை ஒதுக்குவது மட்டுமே
பொருத்தமானது. அணுகுறவுத் தொடர்பு நிலையாக இருந்தால், இயலக்கூடிய வெவ்வேறு
$V(p)$ ஒதுக்கீடுகள்~$W$ இன் எல்லா உட்கணங்களுமே; மேலும்
$\card{W} = n$ எனில் அவை $2^n$ உள்ளன. நமது !!{formula}~$!A$ இல்
$m$ !!{propositional variable}s இருந்தால்,~$n$ உலகங்களைக் கொண்ட
$2^{nm}$ வெவ்வேறு மாதிரிகள் உள்ளன. ஒவ்வொன்றிலும்~$!A$ எல்லா உலகங்களிலும்
மெய்யா என்பதை, ஒவ்வோர் உலகிலும்~$!A$ இன் மெய்மை மதிப்பைக் கணிப்பதன் மூலம்
சோதிக்கலாம். எல்லா இயலக்கூடிய அணுகுறவுத் தொடர்புகளையும் நாம் சோதிக்க
வேண்டும் என்பது உண்மைதான்; ஆனால்~$n$ உலகங்கள் மீதான தொடர்புகளும் முடிவுறு
எண்ணிக்கையிலேயே உள்ளன (குறிப்பாக, $W \times W$ இன் உட்கணங்களின்
எண்ணிக்கை, அதாவது $2^{n^2}$).
% SOURCE CORRECTION TA-NML-030: close the parenthetical counting sentence.

நாம்~\Log{K} தருக்கத்தில் அக்கறை கொள்ளாமல், ஏதோ ஒரு மாதிரி வகுப்பால்
வரையறுக்கப்படும் தருக்கத்தில் (எடுத்துக்காட்டாக, தற்சுட்டும் கடத்தும்
மாதிரிகளில்) அக்கறை கொண்டால், அணுகுறவுத் தொடர்பு சரியான வகையைச் சேர்ந்ததா
என்பதையும் சோதிக்க முடிய வேண்டும். நாம் கருதும் சட்டகங்கள் வாய்ப்புநிலை
!!{formula}s ஆல் வரையறுக்கத்தக்கவையாக இருக்கும் போதெல்லாம் அதைச் செய்யலாம்
(எடுத்துக்காட்டாக, \Ax{T} உம் \Ax{4} உம் அந்தச் சட்டகத்தில் ஏற்புடையவையா
என்று சோதிப்பதன் மூலம்). எனவே, எல்லா முடிவுறு சட்டகங்களையும் ஒன்றன்பின்
ஒன்றாக எடுத்துக் கொண்டு, ஒவ்வொன்றும் நாம் கருதும் வகுப்பிலுள்ள சட்டகமா
என்று சோதித்து, பின்னர் அந்தச் சட்டகத்தின் மீதான எல்லா இயலக்கூடிய
மாதிரிகளையும் பட்டியலிட்டு, ஒவ்வொன்றிலும் $!A$ மெய்யா என்று சோதிப்பதே
கருத்தாகும். மெய்யில்லை என்றால் நிறுத்துக: கருதப்படும் மாதிரி வகுப்பில்
$!A$ ஏற்புடையதல்ல.

இக்கருத்தில் ஒரு சிக்கல் உள்ளது: தேடுவதை எப்போது நிறுத்தலாம், நிறுத்தவே
முடியுமா என்பது நமக்குத் தெரியாது. அந்த !!{formula} க்கு ஒரு முடிவுறு
எதிர்மாதிரி இருந்தால், நமது செயல்முறை அதைக் கண்டுபிடிக்கும். ஆனால் அதற்கு
முடிவுறு எதிர்மாதிரி எதுவும் இல்லையென்றால், நமக்கு விடை கிடைக்காது. அந்த
!!{formula} ஏற்புடையதாக இருக்கலாம் (எதிர்மாதிரிகளே இல்லாமல் இருக்கலாம்),
அல்லது நாம் ஒருபோதும் ஆராயாத ஒரு முடிவிலா எதிர்மாதிரி மட்டுமே அதற்கு
இருக்கலாம். எதிர்மாதிரி உள்ள ஒவ்வொரு !!{formula} க்கும் ஒரு முடிவுறு
எதிர்மாதிரியும் உண்டு என்று காட்ட முடிந்தால் இச்சிக்கலைத் தீர்க்கலாம்.
அப்படி இருப்பின், அந்தத் தருக்கத்துக்கு
\emph{முடிவுறு மாதிரிப் பண்பு} உண்டு என்கிறோம்.

ஆனால் ஒரு தருக்கத்துக்கு முடிவுறு மாதிரிப் பண்பு உண்டு என்பதை எவ்வாறு
காட்டுவது? $!A$ இன் முடிவிலா (எதிர்)மாதிரி ஒன்றை முடிவுறு மாதிரியாக
மாற்றுவதற்கான வழியைக் கண்டுபிடிப்பது ஒரு முறையாகும். அதைச் செய்ய முடிந்தால்,
$!A$ மெய்யாக இல்லாத ஒரு மாதிரி இருக்கும்போதெல்லாம், விளையும் முடிவுறு
மாதிரியிலும் $!A$ மெய்யாகாது. எல்லா முடிவுறு மாதிரிகளின் நமது பட்டியலில்
அந்த முடிவுறு மாதிரி தோன்றும்; இறுதியில், ஏற்புடையதல்லாத ஒவ்வொரு
வாய்பாட்டுக்கும் அது ஏற்புடையதல்ல என்று நாம் தீர்மானிப்போம். வாய்பாடு
ஏற்புடையதாக இருந்தால் நமது செயல்முறை முடிவடையாது. நமது செயல்முறை தரும்
முடிவுறு மாதிரிக்கு ஏதோ ஒரு மிகப்பெரிய அளவு உண்டு என்றும், அந்த மிகப்பெரிய
அளவு !!{formula}~$!A$ ஐ மட்டுமே சாரும் என்றும் கூடுதலாகக் காட்ட முடிந்தால்,
எதிர்மாதிரிகளைத் தேடும்போது முடிவுறு மாதிரிகளை எந்த அளவு வரை சோதிக்க வேண்டும்
என்பது நமக்குக் கிடைக்கும். அதற்குள் ஓர் எதிர்மாதிரியைக் கண்டுபிடிக்கவில்லை
என்றால், எதிர்மாதிரிகள் இல்லை. அப்போது நமது செயல்முறை எந்த
!!{formula}~$!A$ க்கும் ``$!A$ ஏற்புடையதா?'' என்ற கேள்வியை உண்மையிலேயே
தீர்மானிக்கும்.

முடிவிலா அமைப்புகளை முடிவுறு அமைப்புகளாக மாற்றப் பல நேரங்களில் உதவும் ஒரு
உத்தி, பொருத்தமான கூறுகளில் ஒரேபோல் நடந்துகொள்ளும் அமைப்பின் !!{element}s ஐ
``ஒன்றெனக் காண்பது'' ஆகும். பொருத்தமான கூறுகளில் ஒரேபோல் நடந்துகொள்ளும்
முடிவிலா பல உலகங்கள்~$\mModel{M}$ இல் இருந்தால், அத்தகைய உலகங்கள்
\emph{அனைத்தையும்} முடிவுறு எண்ணிக்கையிலான (ஒருவேளை ஒவ்வொன்றும் முடிவிலாவாக
இருக்கக்கூடிய) ``வகுப்புகளாக'' நாம் திரட்ட முடியும். வேறு சொற்களில்,
உலகங்களின் கணத்தைச் சரியான முறையில் கூறிட வேண்டும்; அதாவது, ஒவ்வொரு
கூறிலும் முடிவிலா பல உலகங்கள் இருந்தாலும், கூறுகள் மட்டும் முடிவுறு
எண்ணிக்கையில் இருக்குமாறு கூறிட வேண்டும். பிறகு, கூறுகளையே உலகங்களாகக்
கொண்ட புதிய மாதிரி~$\mModel{M^*}$ ஒன்றை வரையறுக்கிறோம். பழைய மாதிரியிலுள்ள
முடிவுறு எண்ணிக்கையிலான கூறுகள், புதிய மாதிரியில் முடிவுறு எண்ணிக்கையிலான
உலகங்களைத் தருகின்றன; அதாவது, ஒரு முடிவுறு மாதிரியைத் தருகின்றன. உலகம்~$w$
இருக்கும் கூறை $[w]$ என்று அழைப்போம். பழைய மாதிரி~$!A$ க்கு எதிர்மாதிரியாக
இருந்தால் புதிய மாதிரியும் அப்படியே இருக்க வேண்டும் என்பதால்,
$\mSat{M}{!A}[w]$ என்பது $\mSat{M^*}{!A}[{[w]}]$ என்பதற்கு நிகரானதாக
இருக்க வேண்டும் என்று விரும்புகிறோம். இதற்கு, கூறிடலையும்
மாதிரி~$\mModel{M^*}$ இன் அணுகுறவுத் தொடர்பையும் சரியான முறையில் வரையறுக்க
வேண்டும்.

இது எவ்வாறு அமையும் என்பதைப் பார்க்க, முதலில் எல்லா உலகங்களையும் ஒருசேரக்
கருதும் எளிய நிலையை எடுத்துக்கொள்வோம். $\mSat{M}{\Box !B}[w]$ என்பது
ஒவ்வொரு $v \in W$ க்கும் $\mSat{M}{!B}[v]$ என்பதற்கு நிகரானது;
$\mModel{M^*}$ க்கும் இதுவே, ஆனால் $[w]$, $[v]$ வருகின்றன.
% SOURCE CORRECTION TA-NML-031: use the universal Box clause and evaluate B at v.
முதல் கருத்தாக, இரு உலகங்கள் $u$, $v$ ஆகியவை~$\mModel{M}$ இன் எல்லா
!!{propositional variable}s இலும் உடன்பட்டால், அதாவது
$\mSat{M}{p}[u]$ என்பது $\mSat{M}{p}[v]$ என்பதற்கு நிகரானால், அவை சமானமானவை
(ஒரே கூறைச் சேர்ந்தவை) என்று கொள்வோம். $V^*(p) =
\Setabs{[w]}{\mSat{M}{p}[w]}$ என்க. $\mSat{M}{!A}[w]$ என்பது
$\mSat{M^*}{!A}[{[w]}]$ என்பதற்கு நிகரானது என்று காட்டுவதே நமது நோக்கம்.
இதைத் தொகுத்தறிதலால் நிறுவுவோம் என்பது தெளிவு: அடிப்படை நேர்வு
$!A \equiv p$ ஆகும். முதலில் $\mSat{M}{p}[w]$ எனக் கொள்க. அப்போது
வரையறையின்படி $[w] \in V^*$; எனவே $\mSat{M^*}{p}[{[w]}]$. இப்போது
$\mSat{M^*}{p}[{[w]}]$ எனக் கொள்க. இதன் பொருள் $[w] \in V^*(p)$; அதாவது,
~$w$ க்குச் சமானமான ஏதோ ஒரு $v$ க்கு $\mSat{M}{p}[v]$. ஆனால் ``$w$ ஆனது
$v$ க்குச் சமானமானது'' என்பதன் பொருள் ``$w$ உம் $v$ உம் அதே எல்லா
!!{propositional variable}s ஐ மெய்யாக்குகின்றன'' என்பதாகும்; எனவே
$\mSat{M}{p}[w]$. இப்போது தொகுத்தறிதல் படியை, எடுத்துக்காட்டாக,
$!A \ident \lnot !B$ என்பதைக் கருதுக. அப்போது
$\mSat{M}{\lnot !B}[w]$ என்பது $\mSat/{M}{!B}[w]$ என்பதற்கு நிகரானது;
இதுவும் $\mSat/{M^*}{!B}[{[w]}]$ என்பதற்கு (தொகுத்தறிதல்
கருதுகோளின்படி) நிகரானது; இதுவும்
$\mSat{M^*}{\lnot !B}[{[w]}]$. வாய்ப்புநிலையல்லாத மற்ற செய்கைகளுக்கும்
இவ்வாறே. இது $\Box$ க்கும் செயல்படுகிறது: $\mSat{M^*}{\Box
  !B}[{[w]}]$ எனக் கொள்க. இதன் பொருள் ஒவ்வொரு $[u]$ க்கும்
$\mSat{M^*}{!B}[{[u]}]$. தொகுத்தறிதல் கருதுகோளின்படி, ஒவ்வொரு $u$ க்கும்
$\mSat{M}{!B}[u]$. ஆகவே $\mSat{M}{\Box !B}[w]$.

அணுகுறவுத் தொடர்பையும்~$\mModel{M^*}$ க்கு வரையறுக்க வேண்டிய பொதுவான
நேர்வில் நிலை மிகவும் சிக்கலாகும். மேலுள்ள தொகுத்தறிதல் நிறுவலைச் செல்லுபடியாக்க
வேண்டிய நிபந்தனைகளை அதன் அணுகுறவுத் தொடர்பு~$R^*$ நிறைவுசெய்தால்,
மாதிரி~$\mModel{M^*}$ ஒன்றை \emph{வடிகட்டல்} என்று அழைப்போம். அப்போது
எந்த வடிகட்டல்~$\mModel{M^*}$ உம்,~$!A$ ஆனது $[w]$ இல் மெய்யாக இருப்பது
மாதிரி~$\mModel{M}$ ஆனது $!A$ ஐ~$w$ இல் மெய்யாக்குவதற்கு நிகராகுமாறு
செய்யும்.
ஆனால் இப்போது வடிகட்டல்கள் \emph{உள்ளன} என்றும் காட்ட வேண்டும்; அதாவது,
தேவையான நிபந்தனைகளை நிறைவுசெய்யுமாறு~$R^*$ ஐ வரையறுக்க முடியும் என்று
காட்ட வேண்டும். என்றாலும் இது செயல்பட, உலகங்கள் $u$, $v$ ஆகியவை எல்லா
!!{propositional variable}s இலும் உடன்படும்போது மட்டும் அல்லாமல்,~$!A$ இன்
எல்லா துணை-!!{formula}s இலும் உடன்படும்போதுதான் அவை சமானமானவையாகக்
கணக்கிடப்பட வேண்டும். $!A$ க்கு முடிவுறு எண்ணிக்கையிலான துணை-!!{formula}s
மட்டுமே இருப்பதால், வடிகட்டல் முடிவுறுவதாக இருப்பதை இது இன்னும் உறுதிசெய்யும்.
ஒரு வடிகட்டலை வரையறுக்க ஒரே ஒரு வழி மட்டும் இல்லை; வடிகட்டலின் அணுகுறவுத்
தொடர்பு தேவையான பண்புகளை (எடுத்துக்காட்டாக, தற்சுட்டல், கடத்தல் போன்றவற்றை)
நிறைவுசெய்கிறது என்பதை உறுதிப்படுத்த,~$R^*$ இன் வரையறையை நாம்
புத்தாக்கத்துடன் அமைக்க வேண்டும்.


\end{document}
